\subsection{نتایج تجربی}
یافته‌های تجربی حاصل از بنچ‌مارک ماژول \lr{sdrm\_compressor} بر روی مدل‌های زبانی استاندارد در جدول \ref{tab:benchmark_sdrm} منعکس شده است:

\begin{table}[H]
\centering
\small
\caption{نتایج بنچ‌مارک تجربی ماژول فشرده‌ساز مدولاسیون رتبه \lr{sdrm\_compressor}}
\label{tab:benchmark_sdrm}
\begin{tabular}{lccccc}
\toprule
\textbf{مدل پایه \lr{(Base Model)}} & \textbf{دقت \lr{(Accuracy)}} & \textbf{شاخص \lr{F1}} & \textbf{حافظه \lr{(Peak VRAM)}} & \textbf{توان \lr{(Throughput)}} & \textbf{وضعیت} \\
\midrule
\lr{DistilBERT} & \lr{85.78\%} & \lr{0.8631} & \lr{338.98 MB} & \lr{32,743.48 tok/s} & \lr{PASS} \\
\lr{RoBERTa} & \lr{50.92\%} & \lr{0.6748} & \lr{781.79 MB} & \lr{16,739.56 tok/s} & \lr{PASS\_LOW} \\
\lr{GPT-2} & \lr{67.66\%} & \lr{0.7482} & \lr{446.90 MB} & \lr{10,921.04 tok/s} & \lr{PASS} \\
\lr{Qwen} & -- & -- & -- & -- & \lr{SKIP} \\
\bottomrule
\end{tabular}
\end{table}

بر پایه نتایج به دست آمده:
\begin{itemize}
    \item در مدل \lr{DistilBERT}، با فشرده‌سازی فیزیکی ۳۶ لایه و کاهش حجم پارامترها با نرخ فشرده‌سازی حقیقی \lr{1.24x}، دقت پایدار \lr{85.78\%} و شاخص \lr{F1} معادل \lr{0.8631} به دست آمده و وضعیت \lr{PASS} ثبت شده است.
    \item در مدل \lr{GPT-2}، با فشرده‌سازی ۱۵۲ لایه ساختاری و دستیابی به کاهش پارامتر با نسبت \lr{1.26x}، مدل دقت قابل توجه \lr{67.66\%} و شاخص \lr{F1} معادل \lr{0.7482} را در پیش‌بینی توکن‌ها کسب کرده است.
    \item در مدل \lr{RoBERTa}، به دلیل حساسیت شدید لایه‌های توجه این مدل نسبت به فشرده‌سازی فیزیکی ۶۸ لایه بدون فاین‌تیونینگ تکمیلی، دقت مدل به \lr{50.92\%} کاهش یافته که در سطح خطای تصادفی ارزیابی شده و وضعیت \lr{PASS\_LOW} را ایجاد نموده است.
    \item اجرای مدل \lr{Qwen} به علت محدودیت‌های زمانی ارزیابی در این نسخه انجام نگرفته و در وضعیت \lr{SKIP} گزارش شده است.
\end{itemize}
